\section{讨论}
\label{sec:discussion}

\subsection{RQ1：DELB 定义了什么极限？}
\label{subsec:discussion_rq1}

第一个问题是数学问题。对整数值目标，因果信息集下的条件 Shannon 熵可经精确的
整数格点最大熵包络映射为与模型无关的 MSE 底线。预测器相关强化刻画误差仍保留
合法历史信息时的额外代价。普适可达要求误差既与该历史独立，又服从二阶矩匹配的
离散高斯。双松弛恒等式进一步把历史依赖与分布形状失配分开。

更严格地说，比较对象是所有具有有限 MSE 的、可测的整数值因果预测器在给定信息集
下的风险下确界。DELB 不超过这一最优风险。若下确界由某个预测器取得，只有该最优
预测器同时满足两项可达条件时才能写成等号；若最优预测器不存在，即使两个下确界
数值相同，也不能表述为某个预测器达到了 DELB。

这些结论由解析证明得到，不依赖犯罪或自行车应用。数值工作检查配分函数、矩参数化、
逆包络和边界情形的实现，并不是对数学定理的经验验证。相应地，Fano 分析是分类损失
下的经典对照，而不是 DELB 的第二个证明。

\subsection{RQ2：正的 CMI 估计能够说明什么？}
\label{subsec:discussion_rq2}

在总体层面，侧信息使条件熵下降的幅度恰好等于给定基准信息后它与目标的 CMI，
因而不会提高普适 DELB。有限样本估计则是另一个问题。已知真值模拟表明，当总体
CMI 为零时，plugin 与 Miller--Madow CMI 的平均值仍可为正。校正降低了平均
偏差，但在极度稀疏设计中并未一致占优。变换特异的参考分布也可具有正中心，却不必
提高第一类错误率，因为推断比较的是观测统计量与整个变换分布，而不是与零比较。

估计器敏感性分析给出相同警示。在每个城市的 16 个预设规格中，plugin、
Miller--Madow 和基于观测支持的 Jeffreys--Dirichlet 估计均产生正 DELB 降幅。
但 E17 在任何估计器下都没有拒绝：无论在各估计器内部九项城市--零机制检验中控制
错误发现率，还是对全部 27 项估计器--城市--零机制检验进行全局校正，结论均相同。
Jeffreys 结果只是观测支持上的敏感性检查，并不恢复未观测状态。

支持度诊断说明为何必须区分这些对象。加入离散变量会进一步划分基准状态，而随机化
生成的是另一个稀疏状态表。观测原子、singleton 占用率和有效自由度与已知零假设
偏差及经验参考均值均强相关。因此，随机化中心是估计量、样本与变换的共同属性；
它不是总体 CMI，也不是估计量偏差的直接估计。

三轴分析增加了直接的样本量诊断。随着保留更多时间块，相对完整样本DELB和CMI
估计的偏离下降，但在80\%样本时仍不可忽略。该模式描述固定设计下的估计稳定性，
并不表示收集更多观测会改变总体误差底线。同样，面积--时间标准化只是目标单位的
确定性变换，并非不同聚合层级共享同一总体下界的证据。

推断链可写为
\begin{equation}
\begin{gathered}
\text{总体 DELB/CMI estimand}
\longrightarrow
\text{有限样本估计}
\longrightarrow
\text{估计量偏差}\\
\longrightarrow
\text{设计条件随机化参考}
\longrightarrow
\text{实际预测误差}.
\end{gathered}
\label{eq:discussion_inference_chain}
\end{equation}
每个箭头都改变了所回答的问题。若混淆这些层次，就会把一个熵差估计解释成数据并不
支持的更强结论。

\subsection{RQ3：自行车移动是否增加经校准的犯罪预测信息？}
\label{subsec:discussion_rq3}

点估计呈现一致的第一阶段模式：三个城市、全部预设理论替代设定及探索性犯罪类别中，
自行车增强条件熵和 DELB 估计均更低，局部估计也表现出空间结构。这说明增强状态
对观测样本的划分方式降低了估计不确定性。

训练期冻结分析域敏感性保持了这一方向，但改变了幅度。从训练期自行车覆盖域扩展到
训练期完整犯罪域后，纽约和温哥华的 CMI 与 DELB 降幅衰减最明显。因此，估计值
依赖于在验证和测试之前固定的空间目标总体；该分析域对比不具有因果含义。

但校准结果不支持更强的归因。九个城市层面比较和 339 个局部联合筛查均未使观测
统计量与时间、空间和循环移位参考分离。14 日与 30 日移动滞后几乎和一日滞后同样
具有“信息”。因此，对 RQ3 的可辩护回答是：现有设计识别出表面信息增益，但没有
识别出近期、地理匹配且自行车特异的总体 DELB 下降。

嵌套500~m--2~km日与周面板给出相同结论。尽管每个完整样本尺度单元均有正的
DELB下降点估计，270个尺度--样本--零机制比较没有任何一项通过FDR修正。因此，
未获得校准后归因并非仅发生在主要1~km日尺度设定中。

这并不证明条件独立。本文变换是诊断方法，并非在每个基准状态内保持完整移动分布的
精确条件随机化检验。缓慢季节性、疫情期变化、系统运营、天气和遗漏的全市活动都可能
同时存在于近期和远期移动中。该结果限定的是解释强度，而不是证明自行车记录永远不
含有用预测信息。

\subsection{实际预测是另一项后果}
\label{subsec:discussion_prediction}

匹配的留出模型进一步确认式~\eqref{eq:discussion_inference_chain} 中的区分。
若干纽约模型及温哥华状态均值模型在加入自行车状态后改善，而许多 GLM 设定没有
改善。滚动结果也依赖模型与城市。即使熵统计量没有通过自行车特异随机化参考，拟合
模型仍可能利用其中的季节代理信息。

Predictability Gap 仅能作为匹配核算量使用：
\(G=\mathrm{MSE}-L\) 必须采用相同目标、提前期、样本和信息集。更低的估计底线
只创造理论机会，并不保证有限样本估计器能够利用它。Fano 结果在 Hamming 损失下
说明同一点：所有增强底线估计都下降，但透明的众数分类器随测试期未见状态增加而
变差。DELB 与 Fano 针对不同目标和损失，其数值仍不可直接比较。

\subsection{城市数据与预测系统的实际启示}
\label{subsec:practical_implications}

对城市分析团队而言，结果支持分阶段的数据价值评价，而不是依据单个拟合模型做
二元采购判断。表~\ref{tab:practical_decision_framework} 把证据层转换为四个
操作问题。DELB 首先衡量候选数据源是否改变模型无关的估计机会集合；支持度诊断
与随机化随后检验该变化能否在明确的有限样本参考下区分；匹配留出模型最后判断
现有预测系统能否在部署总体中利用该数据源。

\begin{table}[!htbp]
\caption{整数计数预测系统中新增城市数据源的实际决策框架。}
\label{tab:practical_decision_framework}
\centering
\footnotesize
\begin{tabularx}{\textwidth}{p{2.8cm}XXp{3.2cm}}
\toprule
\textbf{决策问题} & \textbf{所需证据} &
\textbf{操作响应} & \textbf{边界}\\
\midrule
数据源是否产生估计预测机会？ &
冻结目标总体中的 DELB 与 CMI 点估计 &
若估计底线实质下降，则保留该数据源进入校准评价 &
点估计不能建立数据源特异信息或实际收益\\
表观增益能否与有限样本结构区分？ &
支持度诊断、已知真值校准与设计匹配随机化 &
仅在校准和功效支持分离时，才对所检验设计进行有限归因 &
低功效下不拒绝没有确定性；随机化不是因果识别\\
现有模型能否在未见数据上利用该数据源？ &
目标、样本、调参和损失均匹配的留出模型 &
只部署表现出稳定留出收益的城市--模型组合，并监测 Predictability Gap &
单一模型改善不能推广到其他城市或模型族\\
价值是否延伸到服务覆盖区外？ &
训练期冻结的服务覆盖总体与完整目标总体 &
同时报告两个总体，并把覆盖范围视为数据产品的一部分 &
服务区估计不能外推到未覆盖地点\\
\bottomrule
\end{tabularx}
\end{table}


这一顺序直接影响数据采购和系统工程。正的 DELB 降幅可以支持保留数据源继续评价，
但不能单独证明应为其付费、完成系统集成或投入运行。本应用中，纽约和温哥华在
自行车覆盖域中的较大降幅扩展到完整犯罪支持总体后明显衰减，校准检验也未分离出
自行车特异信号。因此，覆盖范围、状态支持度和校准应被视为数据产品本身的属性，
而不是事后附加的敏感性检查。

部署判断还应保持城市和模型特异。纽约状态均值与梯度提升模型缩小了
Predictability Gap，其余城市--模型组合没有。实际流水线应保留基准模型，
在冻结留出期评价新增数据源，同时监测
\(\Delta\operatorname{MSE}\) 与 \(\Delta G\)，并在滚动评价中收益不能保持时
撤回新增特征。此时 DELB 是机会预算：它说明估计底线移动多少，而 gap 说明
拟合系统是否把该变化转化为实际精度。

这些启示只针对预测系统评价，不针对实际执法决策。分析是聚合、非因果的，并使用
特定交通方式的移动代理；它既不识别个人，也不指示哪些地点应增加执法，更不识别
自行车活动对犯罪的影响。任何公共部门应用还需要独立的公平、隐私、治理与因果影响
评估，这些均超出 DELB 工作流范围。

\subsection{贡献边界与方法意义}
\label{subsec:discussion_contributions}

本文理论贡献是四部分因果整数预测体系及其与明确推断工作流的连接。离散高斯极值
分布、格点修正、指数族 KL 恒等式和 Fano 不等式均为继承结果。在这一边界内，
本文联合提供：精确逆格点底线、预测器相关强化、可达的充要条件、两种不可达来源的
诊断、侧信息下的总体排序，以及防止把估计底线变化误认为经校准或已实现预测增益的
可复刻有限样本链条。

因此，经验零分布校准不是理论的失败，而是说明模型无关的总体下界在稀疏面板中被
解释前，还需要独立的估计理论。这一区分可迁移到其他整数目标和侧信息来源。

\subsection{局限与下一步}
\label{subsec:limitations}

公共自行车行程是特定交通方式代理，而非 ambient population 的普查。服务边界、
站点布局、采用率、调度、天气、端点缺失和记录规则在城市间不同。2020--2022 年
包含疫情引起的结构变化。报告犯罪受到受害者报案、警方分类与辖区定义的影响，
1~km 日面板也不能支持个体层面的推断。

随机化设计并未在每个基准状态内保持移动的精确条件分布。模拟只覆盖有限的因子设计，
支持度回归是描述性诊断而非解析偏差校正。经验 Fano 底线由测试前分布估计，因此
分布漂移可能使其与测试期总体不等式分离。
基准状态匹配的受限随机化得到相同的不拒绝结果，并在半合成零机制下与名义
第一类错误率相容；但对较小 CMI 的功效仍然有限。因此，该新增实验强化的是
校准审计，而不是“移动与犯罪条件无关”的主张。E18 半合成实验每城仅使用六个
代表网格和 365 日，每单元 50 个 Monte Carlo 数据集及 99 次随机化，二项式
区间较宽；固定经验支持还使若干目标 CMI 不可达。
多尺度随机化在非完整样本下使用一条预设稀释路径以控制计算量；50条路径用于量化
稳定性，但未与每个surrogate完全交叉。空间和时间聚合也会改变目标分布，因此
标准化数值不能识别尺度无关定律。

未来研究可通过收缩、分层汇总或交叉拟合降低状态稀疏性；构造在基准状态内保持移动
条件分布的替代数据；并以匹配因果信息集加入天气、公共交通、行人和 2022 年后的
数据。更尖锐的多元格点包络和替代整数损失仍是理论方向。

\section{结论}
\label{sec:conclusions}

本文沿一条问题链从总体信息论误差底线走向经验证据。离散形式精确的 DELB 为因果
整数预测定义与模型无关的 MSE 下界，并刻画预测器何时能够达到它。总体 CMI 对新增侧信息如何改变
底线进行排序，但稀疏有限样本会使 estimand、点估计和变换特异参考分布彼此分离。
在三个城市面板中，滞后自行车状态降低了 DELB 点估计，但该降幅在主要及嵌套
多尺度设定中均未与所选零机制分离，也没有一致转化为更低的留出误差。核心结论
并非自行车移动缺乏预测价值，而是：
更低的估计底线、经校准的信息增益和改善的实际预测是三种需要不同证据支持的主张。
实际部署新增城市数据源之前，应分别确认估计机会、校准分离和留出收益。
